\documentclass[11pt,a4paper]{article}

% Paquetes
\usepackage[utf8]{inputenc}
\usepackage[T1]{fontenc}
\usepackage[left=25mm,right=25mm,top=25mm,bottom=25mm]{geometry}
\usepackage{graphicx}
\usepackage{multicol}
\usepackage{verbatim}
\usepackage{caption}
\usepackage{subcaption}

\usepackage{pdflscape}
\usepackage{amsmath}
\usepackage[hidelinks]{hyperref}
\usepackage{float}
\usepackage{titlesec}
\usepackage{enumitem}
\usepackage{booktabs}
\usepackage{listings}
\usepackage{xcolor}
\usepackage{tcolorbox}
\usepackage[utf8]{inputenc}
\usepackage[english]{babel}

% Configuración
\captionsetup{font=small}
\titleformat{\section}{\normalfont\Large\bfseries}{\thesection}{1em}{}
\titleformat{\subsection}{\normalfont\large\bfseries}{\thesubsection}{0.75em}{}
\setlist{nosep}
% Listings
\lstset{
  basicstyle=\ttfamily\scriptsize,
  breaklines=true,
  frame=single,
  showstringspaces=false,
  language=python
}

% Nota
\newtcolorbox{nota}{colback=yellow!10,colframe=black!5,boxrule=0.5pt}

\begin{document}

% Portada
 \begin{titlepage}
    \begin{center}
    {\includegraphics[width=1\textwidth]{logo.png}\par}
       \vspace{0.5 cm}
       
    {\bfseries\LARGE Universidad de Zaragoza  \par}
    \vspace{0.5 cm}
    
    {\scshape\Large Facultad de Ingenier\'ia\par}
    \vspace{0.5 cm}
    
    {\scshape\LARGE  Artificial Intelligence - Lab Session 1: Fundamentals\par}
    \vspace{1.5 cm}

    \vfill
    {\Large Jorge Ruiz González (826685) \par}
    {\Large Hugo López Sola (875455) \par}
    
    \vfill
    
    {\Large September 27th, 2026 \par}
    \end{center}
    \end{titlepage}

    \tableofcontents
    \newpage


\section{Aim and evaluation}
We compare different regression and classification methods on
Concrete Compressive Strength, Superconductivity and CIFAR-10.
Scaling and PCA are fitted using training data only. Validation data is used for
model selection, and test data for the final evaluation. RMSE penalizes large
regression errors more than MAE; classification accuracy measures the proportion
of correct predictions, while macro F1 gives equal weight to each class.

\section{Regression}
\subsection{Concrete Compressive Strength}
The dataset contains 1,030 samples and eight input features. We use 721 samples
for training, 206 for validation and 103 for testing. Training correlations with
strength are strongest for \textit{cement} ($r=0.497$), \textit{superplasticizer} ($r=0.377$), \textit{age} ($r=0.328$) and \textit{water} ($r=-0.301$). Zero superplasticizer values are retained:
they can represent a mixture without this additive, rather than an outlier.

Polynomial degree and Ridge regularization are selected by cross-validation (n=5) on the training set. Table~\ref{tab:concrete} compares the models
on validation data.

\begin{table}[ht]
\centering
\small
\caption{Concrete validation results.}
\label{tab:concrete}
\begin{tabular}{lrrr}
\toprule
Model & MAE (MPa) & RMSE (MPa) & $R^2$\\
\midrule
Linear regression & 8.47 & 10.65 & 0.60\\
Polynomial regression, degree 2 & 5.99 & 7.81 & 0.79\\
Polynomial Ridge, degree 3, $\alpha=10$ & 4.95 & 6.51 & 0.85\\
Huber & 8.32 & 10.90 & 0.58\\
RANSAC & 10.32 & 18.65 & $-0.22$\\
\bottomrule
\end{tabular}
\end{table}

Polynomial Ridge achieves the lowest validation RMSE. Its test MAE is
5.20 MPa, RMSE is 7.01 MPa and $R^2=0.798$.
Polynomial features capture nonlinear relationships, while Ridge limits large
coefficients. Huber slightly improves MAE over linear regression but increases
RMSE. RANSAC performs worse on this split; robustness to outliers alone does not
capture nonlinear relationships.

\subsection{Superconductivity}
This dataset contains 21,263 samples and 81 features. The training, validation and
test sets contain 14,884, 4,252 and 2,127 samples. Linear regression obtains the
lowest validation RMSE (17.92 K), although tuned Ridge is practically tied
(17.93 K). PCA retaining 95\% of input variance reduces the data to 17 components,
but PCA followed by Ridge gives an RMSE of 21.99 K. Tuning PCA and Ridge by cross-validation (n=3) improves this to 18.97 K, using 50 components and
$\alpha=10$, still above linear regression.

With a fixed validation set, increasing the training size from 500 to 14,884
reduces linear regression RMSE from 19.21 to 17.92 K. Gains become smaller after
2,000 samples. The selected linear model obtains a test RMSE of
17.44 K and $R^2=0.733$. Preserving input variance with PCA does not
necessarily preserve the information most useful for predicting temperature.

\clearpage
\section{CIFAR-10 classification}
CIFAR-10 contains 60,000 RGB images of $32\times32$ pixels in ten classes.
We select 10,000 training and 2,000 validation images from the official training
set, keeping the official 10,000 test images separate. We compare k-NN ($k=5$),
logistic regression ($C=1$) and SVM with an RBF kernel. SVM parameters $C$ and
$\gamma$ are selected by cross-validation (n=3) on training subsets; the
selected models are then fitted on all 10,000 training images.

For pixel features, scaling and PCA reduce the 3,072 inputs to 64 components.
GIST descriptors are computed for all images using the supplied implementation:
16 Gabor filters and a $4\times4$ spatial grid give 256 features per grayscale
image. Both representations use the same training and validation
images. Table~\ref{tab:cifar} reports test accuracy.

\begin{table}[ht]
\centering
\small
\caption{CIFAR-10 test accuracy for the two representations.}
\label{tab:cifar}
\begin{tabular}{lrr}
\toprule
Model & Pixels + PCA & GIST\\
\midrule
k-NN & 33.5\% & 49.7\%\\
Logistic regression & 37.8\% & 53.8\%\\
SVM & 46.5\% & 62.5\%\\
\bottomrule
\end{tabular}
\end{table}

GIST improves test accuracy by about 16 percentage points for each classifier.
The GIST SVM, selected on validation data, obtains 62.5\% accuracy and
0.625 macro F1. Its per-class F1 is higher for automobile (0.748) and
truck (0.745) than for cat (0.422) and bird (0.494). The notebooks show examples
of correct and incorrect predictions. GIST summarizes texture and spatial
structure, but this grayscale representation loses colour and fine detail.
These results compare the complete feature pipelines, including their different
dimensions and preprocessing.

\section{Conclusions}
Regularized polynomial features improve concrete strength prediction. For
superconductivity, more training data reduces error, while the tested PCA models
do not improve on linear regression. For CIFAR-10, GIST gives higher accuracy
than pixels reduced by PCA, and SVM performs best among the tested classifiers.

\end{document}